\section{Setting and Measurement Levels}
\label{sec:setting}

There are $D$ domains with prompt distributions $\mathcal D_d$, frozen
teachers $q_d$, and a shared student $p_\theta$. A student-generated
response $y$ contains $T$ valid positions with prefixes
$h_t=(x,y_{<t})$. Unless testing a different objective, the fixed-prefix
reference is the reverse-KL loss
\begin{equation}
\ell_d(\theta;h)=\sum_{v\in\mathcal V}
p_\theta(v\mid h)\log\frac{p_\theta(v\mid h)}{q_d(v\mid h)}.
\label{eq:reference_reverse_kl}
\end{equation}
Student and teachers share a tokenizer and vocabulary. Prefixes,
teachers, and loss weights are held fixed when differentiating a local
update. This derivative is not asserted to be the full gradient of a
sequence-level objective whose state distribution depends on $\theta$.

We distinguish the raw gradient $g_d$, the realized optimizer step
$u_d$, the cumulative displacement $\Delta_d=\theta_d-\theta_0$, and
downstream capability $C_d(\theta)$. Sparsity of one object does not
establish sparsity of the others. A vocabulary alternative, a response
position, and a parameter coordinate are also different units.

At a common checkpoint, independently applying each teacher to copies
of the same student makes its local effects comparable. Separately
trained OPD students provide observed single-domain references for
online experiments. They are neither upper bounds nor guarantees of
joint attainability. Capability, loss on a common bank, and loss on
fresh student prefixes are reported separately. A decrease in average
KL alone does not establish balanced capability integration.
